\documentclass[10pt,a4paper]{article}

\usepackage[
    left=1.35cm,
    right=1.35cm,
    top=1.15cm,
    bottom=1.15cm
]{geometry}

\usepackage{enumitem}
\usepackage{hyperref}
\usepackage{xcolor}
\usepackage{titlesec}
\usepackage{fontspec}

\definecolor{darkblue}{HTML}{0F172A}
\definecolor{accentblue}{HTML}{2563EB}

\hypersetup{
    colorlinks=true,
    urlcolor=accentblue,
    linkcolor=darkblue
}

% =========================================================
% GENERAL
% =========================================================

\setlength{\parindent}{0pt}
\setlength{\parskip}{0pt}
\renewcommand{\familydefault}{\sfdefault}
\linespread{1.02}

\raggedbottom

% =========================================================
% SECTION STYLE
% =========================================================

\titleformat{\section}
    {\normalsize\bfseries\color{darkblue}}
    {}
    {0pt}
    {}
    [\titlerule]

\titlespacing*{\section}
    {0pt}
    {0pt}
    {0pt}

% =========================================================
% LISTS
% =========================================================

\setlist[itemize]{
    leftmargin=1.15em,
    itemsep=0pt,
    topsep=0pt,
    parsep=0pt,
    partopsep=0pt
}

\setlist[description]{
    style=multiline,
    leftmargin=3.35cm,
    labelwidth=3.35cm,
    labelsep=0.12cm,
    itemsep=0pt,
    topsep=0pt,
    parsep=0pt,
    partopsep=0pt
}

\begin{document}

% =========================================================
% HEADER
% =========================================================

\begin{center}

    {\LARGE \textbf{ANDRÉ NICOLAS SILVA}} \\[-2pt]

    {\normalsize \textbf{Desenvolvedor de Software | Produtos Web, APIs \& Dados}} \\[1pt]

    \footnotesize
    Brasil \quad • \quad
    \href{mailto:devandrenicolas@gmail.com}{devandrenicolas@gmail.com} \\[-1pt]

    \textbf{\href{https://portfolio-andrenicolas.vercel.app}{Portfólio}}
    \quad • \quad
    \href{https://github.com/devAndreNicolas}{GitHub}
    \quad • \quad
    \href{https://www.linkedin.com/in/devandrenicolas}{LinkedIn}

\end{center}

\vspace{-3pt}

% =========================================================
% SUMMARY
% =========================================================

\section*{RESUMO PROFISSIONAL}

\footnotesize

Desenvolvedor de Software com formação em Sistemas para Internet e experiência prática na construção, manutenção e evolução de produtos SaaS em produção usados por \textbf{10+ empresas}. Atua da compreensão de requisitos e regras de negócio até implementação, testes, depuração e entrega, contribuindo em interfaces web, serviços backend, APIs e dados.

Experiência com \textbf{TypeScript, Angular, React, Go, Node.js, Java/Spring, APIs REST, SQL e PostgreSQL}, além de arquitetura baseada em componentes, integrações, processamento assíncrono, autorização, CI/CD e debugging de produção. Em produto, integrou funcionalidades generativas com Genkit, contexto, sanitização de entradas e respostas estruturadas.

Combina desenvolvimento de software e pensamento de produto para implementar soluções claras e manteráveis. Exerce atualmente o cargo de \textbf{Frontend Engineer na MSPA}, com contribuições full-stack em produtos de privacidade, compliance e auditoria digital.
% =========================================================
% SKILLS
% =========================================================

\section*{COMPETÊNCIAS PRINCIPAIS}

\footnotesize

\begin{description}
    \item[\textbf{Desenvolvimento de Software}]
    Implementação e manutenção de aplicações, resolução de problemas, regras de negócio, arquitetura, code review e suporte à produção

    \item[\textbf{Linguagens \& Frontend}]
    TypeScript, JavaScript, Angular, React, Next.js, HTML/CSS, Angular Signals, RxJS, Web Components, Lit e Shadow DOM

    \item[\textbf{Backend \& Integrações}]
    Go, Node.js, Java e Spring/Spring Boot em contribuições na MSPA; APIs REST/internas, eventos, concorrência, processamento assíncrono e automações

    \item[\textbf{Dados}]
    SQL, PostgreSQL, Supabase, persistência, queries, migrações e validação de regras de negócio

    \item[\textbf{Arquitetura \& Segurança}]
    Hexagonal Architecture, arquitetura baseada em componentes, separação de responsabilidades, autenticação/autorização, validação e auditoria

    \item[\textbf{Qualidade \& Entrega}]
    Vitest, testes automatizados, Testcontainers, Git, CI/CD, debugging, performance e suporte em produção

    \item[\textbf{IA no Produto e Engenharia}]
    Genkit, LLMs, prompt/context engineering, respostas estruturadas, sanitização, harnesses, checks determinísticos e SDD

    \item[\textbf{Cloud \& Containers}]
    Cloudflare Workers, KV, D1, Durable Objects, AWS/Azure em escopo de contribuição, Podman/Docker compatível, Bash
\end{description}

% =========================================================
% EXPERIENCE
% =========================================================

\section*{EXPERIÊNCIA PROFISSIONAL}

\noindent
\textbf{MSPA — Frontend Engineer (contribuições full stack)}
\hfill
\textit{Abr 2025 -- Presente}

\textit{Plataforma SaaS para Privacidade, Compliance (LGPD) e Auditoria Digital}

\begin{itemize}

    \item
    \textbf{Implementação \& Manutenção de Software:}
    Construiu e evoluiu Compass e Trust Center, produtos SaaS usados
    por \textbf{10+ empresas}, implementando funcionalidades a partir
    de requisitos e regras de produto e acompanhando refinamento,
    debugging e entrega em produção.

    \item
    \textbf{Produto \& Tomada de Decisões Técnicas:}
    Traduziu requisitos de negócio e necessidades de produto em
    soluções técnicas, contribuindo para escopo de features,
    regras de negócio, UX, arquitetura, estratégia de
    implementação e trade-offs técnicos.

    \item
    \textbf{Engenharia AI-Native:}
    Aplicou um workflow de desenvolvimento orientado por IA
    utilizando \textbf{Codex, Claude, Cursor, CodeRabbit, Skills,
    MCPs, Agents e SDD} para acelerar implementação, explorar
    soluções, compreender codebases, investigar problemas,
    revisar mudanças e iterar sobre funcionalidades de produto,
    mantendo consistência arquitetural e qualidade de código.

    \item
    \textbf{Engenharia de Produto com IA Generativa:}
    Integração de \textbf{Generative AI em produtos SaaS em
    produção}, desenvolvendo funcionalidades voltadas ao usuário
    com contextual prompt engineering, input sanitization e
    respostas estruturadas de IA utilizando \textbf{Genkit},
    com foco em funcionalidades de produto confiáveis e
    consistentes.

    \item
    \textbf{Engenharia Full Stack:}
    Implementou funcionalidades em produção atravessando
    frontend, backend, APIs, dados e infraestrutura, atuando
    entre diferentes camadas da aplicação de acordo com os
    requisitos de produto e problemas técnicos.

    \item
    \textbf{Stack Ampliada de Plataforma:}
    Contribui na MSPA com \textbf{Node.js, Java, Spring/Spring Boot,
    AWS e Azure/serviços cloud, containers Podman/Docker-compatíveis,
    Bash, Testcontainers, imagens Docker de produção, testes, acessibilidade,
    documentação técnica, entrega de UI e engenharia orientada à segurança},
    selecionando a camada adequada conforme os requisitos do produto.

    \item
    \textbf{Engenharia Backend \& Sistemas:}
    Contribuiu para \textbf{serviços backend e APIs em Golang},
    investigando e implementando soluções envolvendo
    concorrência, processamento assíncrono, eventos, integrações
    entre serviços, automações e problemas de produção em
    diferentes camadas da aplicação.

    \item
    \textbf{Dados \& Regras de Negócio:}
    Trabalhou com \textbf{SQL e PostgreSQL}, incluindo queries,
    persistência de dados, database migrations, regras de negócio
    e investigação de problemas relacionados a dados em
    diferentes workflows da aplicação.

    \item
    \textbf{Arquitetura de Software:}
    Participou de decisões arquiteturais e implementou soluções
    seguindo \textbf{Hexagonal Architecture}, arquitetura baseada
    em componentes, separação de responsabilidades, sistemas
    reutilizáveis e limites de sistema manuteníveis.

    \item
    \textbf{Web \& Edge Systems:}
    Desenvolveu um consent engine agnóstico utilizando
    \textbf{Web Components, Lit e Shadow DOM}, integrado com
    \textbf{Cloudflare Workers, KV, D1 e Durable Objects}
    para execução na edge e funcionalidades distribuídas.

    \item
    \textbf{Performance \& Confiabilidade:}
    Implementou melhorias de performance e SEO, investigando
    gargalos e otimizando comportamento de carregamento e
    performance da aplicação, mantendo usabilidade e
    confiabilidade em produção.

    \item
    \textbf{Entrega de Engenharia:}
    Atuou ao longo do ciclo de entrega de software, incluindo
    \textbf{CI/CD, deployment, debugging, investigação em
    produção e integração de mudanças entre diferentes camadas
    da aplicação}.

    \item
    \textbf{Coordenação de Engenharia:}
    Reorganizou o workflow de desenvolvimento durante uma fase
    crítica do projeto, melhorando clareza de escopo, organização
    de sprints e produtividade de uma \textbf{equipe de 6 pessoas},
    com coordenação direta de 4 integrantes.

\end{itemize}

% =========================================================
% OPEN SOURCE
% =========================================================

\section*{CONTRIBUIÇÕES OPEN SOURCE}

\noindent
\textbf{Stoat / Revolt Platforms LTD — Engenheiro Frontend Voluntário}
\hfill
\textit{2026 -- Presente}

\footnotesize

Contribui para o projeto open source com foco em
\textbf{UX engineering e desenvolvimento frontend}, identificando
problemas voltados aos usuários e transformando-os em melhorias
técnicas que tornam o produto mais claro e intuitivo para os
usuários.
\textbf{Pull request aprovado pelos maintainers do projeto.}

% =========================================================
% PROJECTS
% =========================================================

\section*{PROJETOS}

\noindent
\textbf{Terto Beats — Marketplace e ciclo de pagamento}
\hfill
\textit{Next.js, TypeScript, Supabase, Stripe, Mercado Pago}

\footnotesize

Implementou fluxos de produto digital cobrindo pedidos pendentes, validação de provedor/valor, checkout, webhooks de pagamento e finalização do pedido.

\vspace{0.5pt}

\noindent
\textbf{Marca Racha — Partidas, convites e autorização}
\hfill
\textit{Next.js, React, TypeScript, Supabase, Prisma, Vitest}

\footnotesize

Construiu fluxos para organizar partidas e pedidos de entrada, com autenticação e regras de admissão controladas pelo organizador. Testes comportamentais cobrem permissões, ownership, notificações e transições de estado.

\vspace{0.5pt}

\noindent
\textbf{RendeCerto — Simulador de cenários de investimento}
\hfill
\textit{Next.js, TypeScript, Vitest}

\footnotesize

Implementou regras determinísticas de cálculo e comparação de cenários, com testes automatizados para validar resultados em diferentes entradas.

\vspace{0.5pt}

\noindent
\textbf{Diário de Campo Escoteiro — PWA offline-first}
\hfill
\textit{React, TypeScript, Vite, PouchDB, CouchDB}

\footnotesize

Desenvolveu um diário de campo instalável que mantém registros localmente e sincroniza com CouchDB após a reconexão, usando service worker e manifesto de aplicação web.

\vspace{0.5pt}

\noindent
\textbf{Saúde em Campo — Arquitetura mobile offline-first}
\hfill
\textit{React Native, Expo, SQLite, Drizzle, TanStack Query}

\footnotesize

Desenhou uma aplicação para equipes que registram visitas e informações de saúde em campo mesmo sem conectividade. A proposta combina isolamento local por usuário, fila de sincronização, idempotência, retries/backoff, RBAC e priorização de risco; é evidência de arquitetura, não de implantação clínica.

\vspace{0.5pt}

\noindent
\textbf{Missões do Dia — Aplicativo Android}
\hfill
\textit{Java, SQLite, SharedPreferences, Services, BroadcastReceiver, AsyncTask, XML}

\footnotesize

Desenvolveu um aplicativo acadêmico Android de CRUD de missões com persistência local, componentes nativos e layouts XML.

% EDUCATION
% =========================================================

\section*{FORMAÇÃO}

\noindent
\textbf{Tecnólogo em Sistemas para Internet} — UNCISAL
\hfill
\textit{Concluído em 2026}

% =========================================================
% LANGUAGES
% =========================================================

\vspace{-2pt}
\section*{IDIOMAS}

\noindent
\textbf{Português:} Nativo
\quad • \quad
\textbf{Inglês:} Intermediário/Avançado (em aprimoramento contínuo)

\end{document}
